% ======================================================================
% “华为杯”第二十三届中国研究生数学建模竞赛 - 封面与摘要页模块 (Page 0 & Page 1)
% 
% 说明：
% 1. 本模块严格还原附件3官方 Word 模板的前两页排版规范。
% 2. 依赖宏包（请确保主文档导言区包含）：
%    \usepackage{ctex}
%    \usepackage{graphicx}
%    \usepackage{booktabs}
%    \usepackage{multirow}
%    \usepackage{array}
%    \usepackage[normalem]{ulem}
% 3. 支持在主文档导言区自定义变量宏，若未定义则使用默认/留空值：
%    \newcommand{\SchoolName}{您的学校}
%    \newcommand{\TeamNumber}{您的队号}
%    \newcommand{\MemberA}{队员1}
%    \newcommand{\MemberB}{队员2}
%    \newcommand{\MemberC}{队员3}
%    \newcommand{\PaperTitle}{论文题目}
%    \newcommand{\PaperKeywords}{关键词1\quad 关键词2\quad 关键词3}
%    \newcommand{\AbstractContent}{论文摘要正文...}
% ======================================================================

% -------------------------------------------------------------
% 基础参数保底定义
% -------------------------------------------------------------
\providecommand{\CoverFiguresPath}{figures/} % 徽标图片路径，主文档跨目录引用时可重定义
\providecommand{\SchoolName}{南京理工大学}
\providecommand{\TeamNumber}{}
\providecommand{\MemberA}{}
\providecommand{\MemberB}{}
\providecommand{\MemberC}{}
\providecommand{\PaperTitle}{}
\providecommand{\PaperKeywords}{}
\providecommand{\AbstractContent}{%
（此处填写论文摘要内容。摘要应简明扼要，需包含：建模思路、主要方法、模型、结果与结论、创新点等，请认真书写，注意篇幅一般不超过两页，且无需译成英文。）%
}

% -------------------------------------------------------------
% 字体保底定义（默认使用 ctex 标配中文字体）
% -------------------------------------------------------------
\providecommand{\hwxinwei}{\heiti}
\providecommand{\hwlishu}{\kaishu}

% =============================================================
% 第一页：封面 / 承诺与参赛队伍信息页 (页码标为 0)
% =============================================================
\pagestyle{plain}
\setcounter{page}{0}

% 顶部四校徽/Logo 横向对齐居中
\begin{center}
    \setlength{\tabcolsep}{1.2em}
    \begin{tabular}{@{}cccc@{}}
        \includegraphics[height=1.8cm]{\CoverFiguresPath logo_cpipc.png} &
        \includegraphics[height=2.0cm]{\CoverFiguresPath logo_contest.png} &
        \includegraphics[height=2.1cm]{\CoverFiguresPath logo_huawei.jpg} &
        \includegraphics[height=2.0cm]{\CoverFiguresPath logo_xjtu.png}
    \end{tabular}
\end{center}

\vspace{1.0cm}

% 大赛名称与竞赛名称
\begin{center}
    {\zihao{-2}\hwxinwei\bfseries 中国研究生创新实践系列大赛\par}
    \vspace{0.4cm}
    {\zihao{2}\hwxinwei\bfseries “华为杯”第二十三届中国研究生\par}
    \vspace{0.4cm}
    {\zihao{2}\hwxinwei\bfseries 数学建模竞赛\par}
\end{center}

\vspace{2.2cm}

% 队伍信息表（规范三线表结构）
\begin{center}
\zihao{-2}\bfseries
\renewcommand{\arraystretch}{2.0}
\begin{tabular}{p{3.2cm}p{11.5cm}}
\makebox[4em][s]{学\quad 校} & \SchoolName \\
\specialrule{1.5pt}{0pt}{0pt}
\makebox[4em][s]{参赛队号} & \TeamNumber \\
\specialrule{1.5pt}{0pt}{0pt}
\multirow{3}{*}{\makebox[4em][s]{队员姓名}} & 1.\quad \MemberA \\
\cline{2-2}
 & 2.\quad \MemberB \\
\cline{2-2}
 & 3.\quad \MemberC \\
\specialrule{1.5pt}{0pt}{0pt}
\end{tabular}
\end{center}

\clearpage

% =============================================================
% 第二页：题目与摘要页 (页码从 1 开始)
% =============================================================
\setcounter{page}{1}

% 顶部大赛名称与竞赛名称
\begin{center}
    {\zihao{-2}\hwxinwei\bfseries 中国研究生创新实践系列大赛\par}
    \vspace{0.4cm}
    {\zihao{2}\hwxinwei\bfseries “华为杯”第二十三届中国研究生\par}
    \vspace{0.4cm}
    {\zihao{2}\hwxinwei\bfseries 数学建模竞赛\par}
\end{center}

\vspace{1.0cm}

% 题目栏
\noindent{\zihao{-2}\hwlishu 题\quad 目：}\quad 
\ifx\PaperTitle\empty
    \uline{\hspace{11.5cm}}
\else
    \uline{\makebox[11.5cm][c]{\zihao{3}\songti\bfseries \PaperTitle}}
\fi
\par

\vspace{1.2cm}

% 摘要标题（居中）
\begin{center}
    {\zihao{-2}\hwlishu 摘\quad 要：\par}
\end{center}

\vspace{0.5cm}

% 摘要正文区域（小四宋体，1.3倍行距）
\begingroup
\zihao{-4}\songti
\linespread{1.3}\selectfont
\AbstractContent
\endgroup

\vfill

% 关键词栏
\noindent{\zihao{-2}\hwlishu 关键词：}\quad 
\begingroup
\zihao{-4}\songti
\ifx\PaperKeywords\empty
    \quad
\else
    \PaperKeywords
\fi
\endgroup
\par

\vspace{0.8cm}
\clearpage
